\chapter{商空间}

\section{商拓扑与商空间}

\begin{definition}[商空间]
    设 $(X, \mathcal{T})$ 是一个拓扑空间，$\sim$ 是 $X$ 上的一个等价关系.
    \begin{itemize}
        \item $X/\!\sim$ 是商集.
        \item $\setjot[-2]\begin{aligned}[t]
            q: X &\to X/\!\sim \\
            x &\mapsto [x]
        \end{aligned}$ 是典范映射.
    \end{itemize}
    令 $\tilde{\mathcal{T}} := \{ U \subset X/\!\sim \mid q^{-1}(U) \in \mathcal{T} \}$.则$\tilde{\mathcal{T}}$ 是 $X/\!\sim$ 上的一个拓扑，称为\textbf{商拓扑}. $(X/\!\sim, \tilde{\mathcal{T}})$ 称为\textbf{商空间}，$q:X\to X/\!\sim$ 称为\textbf{商映射}.
\end{definition}
\begin{proof}[证明“$\tilde{\mathcal{T}}$是一个拓扑”]
    (i) $\varnothing, X/\!\sim \in \tilde{\mathcal{T}}$，因为 $q^{-1}(\varnothing)=\varnothing \in \mathcal{T}$，$q^{-1}(X/\!\sim)=X \in \mathcal{T}$.

    (ii) 设 $U_\lambda \in \tilde{\mathcal{T}}, \lambda \in \Lambda$.则 $q^{-1}(\bigcup_{\lambda \in \Lambda} U_\lambda) = \bigcup_{\lambda \in \Lambda} q^{-1}(U_\lambda)$.由于每个 $q^{-1}(U_\lambda)$ 都是开集，它们的并集也是开集，因此 $\bigcup_{\lambda \in \Lambda} U_\lambda \in \tilde{\mathcal{T}}$.

    (iii) 设 $U_1, U_2 \in \tilde{\mathcal{T}}$.则 $q^{-1}(U_1 \cap U_2) = q^{-1}(U_1) \cap q^{-1}(U_2)$.由于 $q^{-1}(U_1)$ 和 $q^{-1}(U_2)$ 都是开集，它们的交集也是开集，因此 $U_1 \cap U_2 \in \tilde{\mathcal{T}}$.
\end{proof}

\begin{remark}
    (1) 商映射 $q: X \to X/\!\sim$ 是满射.\par
    (2) $U \subset X/\!\sim$ 是开集 $\iff q^{-1}(U) \subset X$ 是开集.
\end{remark}

\begin{instance}
    令 $q: \mathbb{E}^1 \to \{a, b, c\}$ 定义为
    \[ q(x) = \begin{cases} a, & x > 0 \\ b, & x = 0 \\ c, & x < 0 \end{cases} \]
    则 $\{a,b,c\}$ 上的商拓扑为 $\{\varnothing, \{a\}, \{c\}, \{a, c\}, \{a, b, c\}\}$.
\end{instance}

\begin{instance}
    若 $(X, \mathcal{T})$ 是离散拓扑，则对于 $X$ 上的任意等价关系 $\sim$，$(X/\!\sim, \tilde{\mathcal{T}})$ 也是离散拓扑.
\end{instance}

\begin{proposition}
    商拓扑是使得 $q: X \to X/\!\sim$ 连续的最精细的拓扑.
\end{proposition}
\begin{proof}
    假设 $\tilde{\mathcal{T}}'$ 是 $X/\!\sim$ 上的一个拓扑，且 $q: (X, \mathcal{T}) \to (X/\!\sim, \tilde{\mathcal{T}}')$ 连续.
    那么对于任意 $U \in \tilde{\mathcal{T}}'$，$q^{-1}(U) \in \mathcal{T}$.由商拓扑的定义，这意味着 $U \in \tilde{\mathcal{T}}$.
    因此 $\tilde{\mathcal{T}}' \subset \tilde{\mathcal{T}}$.
\end{proof}

\begin{lemma}\label{lem:quotient_map_char}
    设 $q: X \to X/\!\sim$ 是一个商映射，$Z$ 是一个拓扑空间，$f: X/\!\sim \to Z$ 是一个映射.
    则 $f$ 连续 $\iff f \circ q$ 连续.
\end{lemma}
\begin{proof}
    “$\Rightarrow$” 两个连续映射的复合是连续的，显然成立.

    “$\Leftarrow$” 设 $V$ 是 $Z$ 中的任意开集.我们需要证明 $f^{-1}(V)$ 在 $X/\!\sim$ 中是开集.
    根据商拓扑的定义，这等价于证明 $q^{-1}(f^{-1}(V))$ 在 $X$ 中是开集.
    注意到 $q^{-1}(f^{-1}(V)) = (f \circ q)^{-1}(V)$.
    因为 $f \circ q$ 是连续的，所以 $(f \circ q)^{-1}(V)$ 在 $X$ 中是开集.
    因此 $f^{-1}(V)$ 在 $X/\!\sim$ 中是开集，故 $f$ 连续.
    \[
    \begin{tikzcd}
        X \arrow[rd, "f \circ q"] \arrow[d, "q"'] & \\
        X/\!\sim \arrow[r, "f"] & Z
    \end{tikzcd}
    \]
\end{proof}

\begin{lemma}
    设 $q: X \to X/\!\sim$ 是一个商映射.则 $C \subset X/\!\sim$ 是闭集 $\iff q^{-1}(C)$ 在 $X$ 中是闭集.
\end{lemma}

\begin{definition}[商映射的抽象定义]
    设 $X, Y$ 是拓扑空间.如果映射$f: X \to Y$满足：
    \begin{enumerate}[(1)]
        \item $f$ 是满射.
        \item $U \subset Y$ 是开集(闭集) $\iff f^{-1}(U) \subset X$ 是开集(闭集).
    \end{enumerate}
    则称$f$为一个\textbf{商映射}.
\end{definition}

\begin{remark}
    引理 \ref{lem:quotient_map_char} 对于抽象定义的商映射也成立.
\end{remark}

\begin{theorem}[商空间同胚定理]
    设 $f: X \to Y$ 是一个商映射，$x_1, x_2 \in X$.定义等价关系 $x_1 \sim_f x_2 \iff f(x_1) = f(x_2)$.
    则 $X/\!\sim_f \cong Y$.
\end{theorem}
\begin{proof}
    定义 $h: X/\!\sim_f \to Y, [x] \mapsto f(x)$.容易验证 $h$ 是良定义且是双射.
    我们需要证明 $h$ 和 $h^{-1}$ 都是连续的.
    \begin{itemize}
        \item 证明 $h$ 连续：
        考虑交换图表：
        \[
        \begin{tikzcd}
            X \arrow[rd, "f"] \arrow[d, "q"'] & \\
            X/\!\sim_f \arrow[r, "h"] & Y
        \end{tikzcd}
        \]
        我们有 $h \circ q = f$.因为 $f$ 是连续的，根据引理 \ref{lem:quotient_map_char}，$h$ 是连续的.

        \item 证明 $h^{-1}$ 连续：
        考虑交换图表：
        \[
        \begin{tikzcd}
            X \arrow[d, "f"'] \arrow[r, "q"] & X/\!\sim_f \\
            Y \arrow[ru, "h^{-1}"'] &
        \end{tikzcd}
        \]
        我们有 $h^{-1} \circ f = q$.因为 $q$ 是连续的，而 $f$ 是商映射（根据抽象定义），再次使用引理 \ref{lem:quotient_map_char} 的变体（将 $q$ 换成 $f$），可知 $h^{-1}$ 是连续的.
    \end{itemize}
    因此，$h$ 是一个同胚.
\end{proof}

\begin{proposition}
    若映射 $f:X\to Y$ 是满射、连续、开（或闭）映射，则 $f$ 是一个商映射.
\end{proposition}
\begin{proof}
    假设 $U\subset Y$ 且 $f^{-1}(U)\subset X$ 是开集.由于 $f$ 是开映射，所以 $f(f^{-1}(U))$ 在 $Y$ 中是开集.又因为 $f$ 是满射，所以 $f(f^{-1}(U))=U$.因此 $U$ 在 $Y$ 中是开集.根据商映射的定义，$f$ 是一个商映射.
    \par
    若 $f$ 是闭映射，证明是类似的.
\end{proof}

\begin{instance}
    投影映射 $\pi_X:X\times Y\to X, (x,y)\mapsto x$ 是一个商映射，因为 $\pi_X$ 是满射、连续且开的.
\end{instance}

\begin{proposition}
    设 $f:X\to Y$ 是一个满射、连续的映射.若 $X$ 是紧空间，$Y$ 是Hausdorff空间，则 $f$ 是一个商映射.
\end{proposition}
\begin{proof}
    只需证明 $f$ 是一个闭映射即可.
    \par
    设 $C$ 是 $X$ 中的一个闭集.由于 $X$ 是紧空间，所以 $C$ 是 $X$ 中的紧子集.
    \par
    由于 $f$ 连续，所以 $f(C)$ 是 $Y$ 中的紧子集.
    \par
    由于 $Y$ 是Hausdorff空间，所以 $f(C)$ 在 $Y$ 中是闭集.
\end{proof}

\begin{lemma}
    设 $f:X\to Y$ 和 $g:Y\to Z$ 都是商映射，则复合映射 $g\circ f:X\to Z$ 也是一个商映射.
\end{lemma}
\begin{proof}
    $g\circ f$ 是满射和连续的.假设 $U\subset Z$ 且 $(g\circ f)^{-1}(U)$ 在 $X$ 中是开集.
    \par
    $(g\circ f)^{-1}(U)=f^{-1}(g^{-1}(U))$.
    \par
    由于 $f$ 是商映射，因此 $g^{-1}(U)$ 在 $Y$ 中是开集.
    \par
    由于 $g$ 是商映射，因此 $U$ 在 $Z$ 中是开集.
\end{proof}

\begin{remark}
    (1) 设 $f_1:X_1\to Y_1,f_2:X_2\to Y_2$ 是商映射，则 $f_1\times f_2:X_1\times X_2\to Y_1\times Y_2$ 不一定是商映射.
    \par
    (2) 设 $f:X\to Y$ 是商映射.若 $X$ 是Hausdorff空间，$Y$ 不一定是Hausdorff空间.
    \par
    (3) 商映射不一定是开映射或闭映射.
\end{remark}

\section{例子}

\begin{instance}
    设 $X$ 是一个拓扑空间，$A\subset X$.定义等价关系 $x_1\sim x_2 \iff x_1,x_2\in A$ (对所有 $x_1\neq x_2$ 而言).商空间 $X/\!\sim$ 记作 $X/A$.这相当于把子集 $A$ “压缩”成一个点.
\end{instance}

\begin{instance}
    设 $D^n$ 是 $n$ 维闭圆盘，$S^{n-1}\subset D^n$ 是其边界（一个 $n-1$ 维球面），$n\geq 1$.则 $D^n/S^{n-1}\cong S^n$.
    \par
    考虑映射 $f:D^n\to S^n\subset \mathbb{E}^{n+1}$，定义为
    \[ (r,s)\mapsto (2\sqrt{r(1-r)}s, 2r-1),\quad 0\leq r\leq 1, s\in S^{n-1} \]
    这里我们把 $D^n$ 的点用极坐标类似的方式表示. $f$ 是满射、连续的.由于 $D^n$ 是紧的，$S^n$ 是Hausdorff的，所以 $f$ 是一个商映射.
    \par
    因此 $D^n/S^{n-1}=D^n/\!\sim_f\cong S^n$.
\end{instance}

\begin{instance}
    设 $X$ 是一个拓扑空间.将 $X\times [0,1]$ 中的子空间 $X\times\{1\}$ 压缩成一个点，得到的商空间 $X\times[0,1]/X\times\{1\}$ 称为 $X$ 的\textbf{拓扑锥(topological cone)}，记作 $CX$.
\end{instance}

\begin{instance}[一些由正方形粘合得到的曲面]
    \begin{center}
        \includegraphics[width=0.8\textwidth]{example-image}
    \end{center}
    \begin{itemize}
        \item 圆环 (Annulus)
        \item 莫比乌斯带 (Möbius band)
        \item 环面 (Torus)
        \item 克莱因瓶 (Klein bottle)
        \item 球面 (Sphere)
        \item 射影平面 (Projective plane)
    \end{itemize}
\end{instance}

\begin{instance}[射影平面的构造]
    \begin{center}
        \includegraphics[width=\textwidth]{example-image}
    \end{center}
    上图展示了通过粘合一个圆盘的边界上的对径点来构造射影平面的过程.
\end{instance}

\begin{instance}[曲面剪贴]
    \begin{itemize}
        \item 圆环(Annulus) $\cong$ 莫比乌斯带(Möbius band)
        \begin{center}
            \includegraphics[width=0.8\textwidth]{example-image}
        \end{center}
        上图将圆柱体（拓扑上是圆环）剪开、扭转其中一条边再重新粘合，从而得到一个莫比乌斯带.
        \item 交叉帽(Cross cap)
        \begin{center}
            \includegraphics[width=0.6\textwidth]{example-image}
        \end{center}
        交叉帽是射影平面的另一种嵌入表示.
        \item 莫比乌斯带 $\cup$ 圆盘 $\cong$ 射影平面($\mathbb{P}^2$)
        \begin{center}
            \includegraphics[width=0.8\textwidth]{example-image}
        \end{center}
        将一个圆盘的边界沿着莫比乌斯带的边界粘合，可以得到一个射影平面.
        \item 三角形 $\cong$ 莫比乌斯带
        \begin{center}
            \includegraphics[width=0.6\textwidth]{example-image}
        \end{center}
        \item 莫比乌斯带 $\cup$ 莫比乌斯带 $\to$ 克莱因瓶(Klein bottle)
        \begin{center}
            \includegraphics[width=\textwidth]{example-image}
        \end{center}
        将两个莫比乌斯带沿着它们的边界粘合在一起，可以得到一个克莱因瓶.
    \end{itemize}
\end{instance}

\section{射影空间与轨道空间}

\begin{definition}[射影空间]
    $n$维\textbf{射影空间}(projective space)，记作 $\mathbb{P}^n$，有以下几种等价的构造方式：
    \begin{enumerate}[(i)]
        \item 在 $\mathbb{E}^{n+1}\setminus\{0\}$ 中，定义等价关系：对于 $\vec{y}_1, \vec{y}_2 \in \mathbb{E}^{n+1}\setminus\{0\}$，
        \[ \vec{y}_1 \sim \vec{y}_2 \iff \exists \lambda \in \mathbb{R}\setminus\{0\}, \st \vec{y}_2 = \lambda \vec{y}_1 \]
        （即，所有过原点的直线上的非零点都等价）.
        商空间 $(\mathbb{E}^{n+1}\setminus\{0\})/\sim$ 定义为 $\mathbb{P}^n$.
        \item 在 $n$维球面 $S^n = \{ (x_0, \dots, x_n) \in \mathbb{E}^{n+1} \mid x_0^2 + \dots + x_n^2 = 1 \}$ 中，定义等价关系：对于 $\vec{y}_1, \vec{y}_2 \in S^n$，
        \[ \vec{y}_1 \sim \vec{y}_2 \iff \vec{y}_2 = -\vec{y}_1 \]
        （即，将对径点视为同一点）.
        则 $S^n/\sim \cong \mathbb{P}^n$.
        \item 在 $n$维闭半球面 $D^n = \{ (x_0, \dots, x_n) \in S^n \mid x_n \geq 0 \}$ 中，定义等价关系：对于 $\vec{y}_1, \vec{y}_2 \in D^n$，
        \[ \vec{y}_1 \sim \vec{y}_2 \iff \vec{y}_1, \vec{y}_2 \in \partial D^n = S^{n-1} \text{ 且 } \vec{y}_2 = -\vec{y}_1 \]
        （即，只将边界（赤道）上的对径点视为同一点）.
        则 $D^n/\sim \cong \mathbb{P}^n$.
    \end{enumerate}
\end{definition}

\begin{definition}[轨道空间]
    设 $(X, \mathcal{T})$ 是一个拓扑空间， $G$ 是一个群.若 $G$ 在 $X$ 上有群作用（$G$ acts on $X$），则可以定义等价关系：对于 $x_1, x_2 \in X$，
    \[ x_1 \sim x_2 \iff \exists g \in G, \st x_2 = g \cdot x_1 \]
    商空间 $X/\sim$ 称为\textbf{轨道空间}(orbit space)，记作 $X/G$.
\end{definition}

\begin{instance}[轨道空间的例子]
    \begin{enumerate}[(i)]
        \item 整数群 $\mathbb{Z}$ 通过平移作用于实直线 $\mathbb{E}^1$：$z \in \mathbb{Z}$ 的作用为 $x \mapsto x+z$.
        轨道空间 $\mathbb{E}^1/\mathbb{Z} \cong S^1$. 这相当于将直线卷成一个圆.
        \item 群 $\mathbb{Z}\oplus\mathbb{Z}$ 通过平移作用于平面 $\mathbb{E}^2$：$(z_1, z_2) \in \mathbb{Z}\oplus\mathbb{Z}$ 的作用为 $(x,y) \mapsto (x+z_1, y+z_2)$.
        轨道空间 $\mathbb{E}^2/(\mathbb{Z}\oplus\mathbb{Z}) \cong T^2$. 这相当于将平面卷成一个环面.
        \item 2阶循环群 $\mathbb{Z}_2 = \{+1, -1\}$ 通过对径映射作用于球面 $S^n$：$-1$ 的作用为 $x \mapsto -x$.
        轨道空间 $S^n/\mathbb{Z}_2 \cong \mathbb{P}^n$. 这正是射影空间的构造(ii).
    \end{enumerate}
\end{instance}

\begin{definition}[透镜空间]
    设 $p, q$ 是整数，满足 $p>2, 1\leq q < p, \gcd(p,q)=1$.
    令 $g$ 为循环群 $\mathbb{Z}_p$ 的一个生成元. $\mathbb{Z}_p$ 作用在三维球面 $S^3 = \{ (z_1, z_2) \in \mathbb{C}^2 \mid |z_1|^2+|z_2|^2=1 \}$ 上，其作用方式为：
    \[ g \cdot (z_1, z_2) = \left(z_1 e^{2\pi i/p}, z_2 e^{2\pi i q/p}\right) \]
    所得到的轨道空间 $S^3/\mathbb{Z}_p$ 称为\textbf{透镜空间}(lens space)，记作 $L(p,q)$.
\end{definition}




